\documentclass[10pt,a4paper]{amsart}
\usepackage[margin=19mm]{geometry}
\usepackage{amsmath,amssymb,mathtools}
\usepackage[scr=rsfs,cal=euler]{mathalfa}
\usepackage{xfrac}
\usepackage[unicode,hidelinks,bookmarks=false]{hyperref}
\newcommand{\ie}{i.e.\ }
\newcommand{\cf}{cf.\ }
\newcommand{\resp}{respectively\ }
\newcommand{\Subsection}{\subsection}
\providecommand{\nobreakdash}{\nobreak\hbox{-}}
\setlength{\emergencystretch}{2em}
\multlinegap0pt
\usepackage{amssymb}
\newtheorem{proposition}{Proposition}
\title{On Discrete Approximations to Infinite Horizon Differential Games}
\author{Javier de Frutos, V\'ictor Gat\'on, Julia Novo}
\date{Source: arXiv:2112.03153v3 (CC BY 4.0)}
\begin{document}
\maketitle

\noindent\emph{Source and licence.} On Discrete Approximations to Infinite Horizon Differential Games. Version arXiv:2112.03153v3 (CC BY 4.0), distributed by arXiv under the Creative Commons Attribution 4.0 International licence, http://creativecommons.org/licenses/by/4.0/, as recorded in the arXiv licence metadata for this exact version and shown on the abstract page https://arxiv.org/abs/2112.03153v3. Full licence notice: \texttt{licenses/arxiv-2112.03153-CC-BY-4.0.txt}.\par\smallskip

\noindent\textit{Editorial scope.} Counted unit: the article's proposition consitency\_with\_order (source0.tex lines 320--327). The article's complete original proof of it is delivered with it (source0.tex lines 328--385, one \texttt{proof} environment of 58 lines). No mathematics is written, repaired or re-derived: the statement and the proof are the article's own lines, in the article's own notation, and no retained line is substituted or re-flowed.  Retained uncounted as the source-local context the proof needs: lines 206--217; lines 246--248; lines 264--266.  Nothing was omitted from any retained span, so the package carries no omission record.  No retained line contains a citation, so the wrapper carries no bibliography and no bibliography entry.  No retained line contains a figure environment, an image-inclusion command or drawing code, so no image is delivered and the package declares no asset dependency.  Editorial formatting only, in addition: an AMS \texttt{amsart} preamble that restates the article's own declarations and macros used by the retained lines, namely the environments \texttt{proposition} on separate counters, together with the standard packages those lines need.  The retained environments print in this document as Proposition 1 in order of appearance, whereas the article prints the same items with the numbering of its own full document.  The complete native source of the article as distributed in the arXiv e-print for this exact version is bundled unchanged as \texttt{sources/119327/source0.tex}.
\begin{proposition}\label{consistency}
Let  $(\phi_1,\dots,\phi_N)\in\mathcal{U}$ an arbitrary $N$-tuple of admissible strategies. Let us assume that there exists a constant $L_s>0$ with
\begin{equation}\label{lipschitz_strategies}
|\phi_i(x_1)-\phi_i(x_2)|\le L_s |x_1-x_2|,\quad i=1,\dots, N.
\end{equation}
Let us assume that hypotheses $H_1$, $H_2$ and $H_3$ are satisfied.
 Then
\begin{equation*}
\lim_{h\rightarrow 0} \left|W_{i,h}(\phi_{i},\phi_{-i},x)- W_i(\phi_{i},\phi_{-i},x)\right|=0,
\quad i=1,\dots,N.
\end{equation*}
\end{proposition}

\begin{equation}\label{bound_dynamics}
|g(y,u_{i},u_{-i})|\le K\big(1+|y|\bigr),
\end{equation}

\begin{equation}\label{estimacion_dynamics}
|y(t)-\tilde y(t)|\le K h\left(1+(|x|+\sqrt{2Kt})e^{Kt}\right)e^{Lt}.
\end{equation}

\begin{proposition}\label{consitency_with_order}
Let  $(\phi_1,\dots,\phi_N)\in\mathcal{U}$ an arbitrary $N$-tuple of {(time-continuous)} admissible strategies and let $x\in\mathbb{V}$. Let assume that hypotheses $H1$, $H2$, $H3$ and \eqref{lipschitz_strategies} hold.
Let us assume that either {$\rho> L$} with $L=L_g(1+LN_s)$ and $H4$ holds or {$\rho>L+K$} with $K$ the constant in \eqref{bound_dynamics}.  Then, there exists a
positive constant $C$ and $h_0>0$ such that for all $h\le h_0$ {with $h_0<1/\rho$}
\begin{equation*}
\bigl|W_{i,h}(\phi_i,\phi_{-i},x)-W_i(\phi_i,\phi_{-i},x)\bigr |\le C h.
\end{equation*}
\end{proposition}

\begin{proof}
In the proof we will use the same notation as in Proposition \ref{consistency}.

We start by writing
\begin{align} \label{auxiliar_consitency_with_order}
\Bigl |W_{i,h}({\phi}_i,{\phi}_{-i}&,x)- W_i(\phi_i,\phi_{-i},x)\Bigr |\le\nonumber\\
&\int_0^\infty |f_i(y(t),\phi_i(y(t)),\phi_{-i}(y(t)))-f_i(\tilde{y}(t),{\phi}_i(\tilde{y}(t)),\phi_{-i}(\tilde{y}(t)))|e^{-\rho t}\, ds,\nonumber\\
&+\int_0^\infty|f_i(\tilde{y}(t),{\phi}_i(\tilde{y}(t)),\phi_{-i}(\tilde{y}(t)))||e^{-\rho t}-e^{- \rho\theta_h[t/h]h}|\, dt,
\end{align}
where $\theta_h=-\log(1-\rho h)/(\rho h)$.
Then, using  H2, (\ref{lipschitz_strategies}) and (\ref{estimacion_dynamics}) we have, if {$\rho>K+ L$},
\begin{equation*}
\begin{split}
\int_0^\infty |f_i(y(t),\phi_i(y(t)),&\phi_{-i}(y(t)))-f_i(\tilde{y}(t),{\phi}_i(\tilde{y}(t)),\phi_{-i}(\tilde{y}(t)))|e^{-\rho t}\, dt\le \\
&hKL_i(1+NL_s)\int_0^\infty\left(1+(|x|+\sqrt{2Kt})e^{Kt}\right)e^{(L-\rho)t} \, dt \le C h,
\end{split}
\end{equation*}
for some constant $C>0$.

The second term in (\ref{auxiliar_consitency_with_order}) can be estimated using hypothesis H3 and the mean
 value theorem {applied to the function $e^{-\rho s}$}

\begin{equation*}
\begin{split}
\int_0^\infty|f_i(\tilde{y}(t),{\phi}_i(\tilde{y}(t)),\phi_{-i}(\tilde{y}(t)))|&|e^{-\rho t}-e^{- \rho\theta_h[t/h]h}|\, dt\le\\&{
M_f \int_0^\infty \rho\max\left\{e^{-\rho t},e^{-\rho\theta_h[t/h]h}\right\}\rho| t- \theta[t/h]h| \, dt.}
\end{split}
\end{equation*}
{Taking into account that $[t/h]h\le t\le [t/h]h+h$ and $\theta_h>1$ then $|t- \theta[t/h]h|\le (\theta_h-1)t+\theta_h h$. On the other hand, it is easy to check that $\max\left\{e^{-\rho t},e^{-\rho\theta_h[t/h]h}\right\}
\le e^{-\rho t} e^{\rho \theta_h}.$ Then}
\begin{equation*}
\begin{split}
\int_0^\infty|f_i(\tilde{y}(t),{\phi}_i(\tilde{y}(t)),\phi_{-i}(\tilde{y}(t)))|&|e^{-\rho t}-e^{- \rho\theta_h[t/h]h}|\, dt\le\\
&M_f\rho^2 {e^{\theta_h\rho }}\int_0^{\infty}e^{-\rho t}((\theta_h-1)t+\theta_h h) \, dt,\\
\end{split}
\end{equation*}
and since
 $ \theta_h-1=\mathcal{O}(h)$ as $h\rightarrow 0$, we conclude
\begin{equation*}
\int_0^\infty|f_i(\tilde{y}(t),{\phi}_i(\tilde{y}(t)),\phi_{-i}(\tilde{y}(t)))||e^{-\rho t}-e^{- \rho\theta_h[t/h]h}|\, dt\le Ch,
\end{equation*}
for some constant $C>0.$

Let us assume $H4$ and {$\rho>L$}.
We have immediately that{
$$
|y(t)-\tilde y(t)|\le M_g h e^{Lt},
$$
}
and  then
\begin{equation*}
\begin{split}
\int_0^\infty |f_i(y(t),\phi_i(y(t)),\phi_{-i}(y(t)))-f_i(\tilde{y}(t)&,{\phi}_i(\tilde{y}(t)),\phi_{-i}(\tilde{y}(t)))|e^{-\rho t}\, ds\le \\
&L_i(1+NL_s)M_gh\int_0^\infty e^{(L-\rho)t}\, dt \le Ch,
\end{split}
\end{equation*}
for some positive constant $C>0$ which finishes the proof.
\end{proof}

\end{document}
